\subsection{Simplifying Roots of Numbers} 
Often, when we solve an equation using the root properties, we are left with an irrational root in our solution.
\begin{example}
Solve $x^3=54$ using the appropriate root property.
\begin{lpic}{}{2}
\end{lpic}
\end{example}
\noi We know we can get a decimal approximation for this number. We can also \term{simplify} it. Recall:
\begin{definition}
To \term{simplify} an expression is to write an equivalent expression in some standard ``simplest'' from.
\end{definition}
To simplify a radical expression, we use the following property of roots (which is the same as a property of powers).
\begin{theorem}[Root of a Product Rule]
For any real numbers $x$ and $y$ and any positive integer $p$ where all the roots below are real numbers,

\nic$\sqrt[p]{xy}=\sqrt[p]{x}\cdot\sqrt[p]{y}$
\end{theorem}
To simplify a root, then, is to take any possible integer roots.
\begin{definition}
For any real number $x$ and integer index $p$, the radical expression $\sqrt[p]{x}$ is said to be \term{simplified} if it is written in the form $a\sqrt[p]{b}$ where $a^pb=x$ and $b$ is as small as possible.
\end{definition}
To accomplish this, we need to find the largest possible factor of $x$ that is a perfect $p$th power.
\begin{example}Simplify $\sqrt[3]{54}$.
\begin{lpic}{}{4}
\foreach \y in {1,...,4}
	\regtextleft{0.25}{-\y}{$\y^3={}$};
\foreach \y in {1,...,3}
	\regtextleft{4.25}{-\y}{$54\div{}\qquad={}$};
\regtextleft{10.25}{-1}{$\sqrt[3]{54}={}$};
\end{lpic}
\end{example}